\documentclass[10pt,a4paper]{amsart}
\usepackage[margin=19mm]{geometry}
\usepackage{amsmath,amssymb,mathtools}
\usepackage[scr=rsfs,cal=euler]{mathalfa}
\usepackage{xfrac}
\usepackage[unicode,hidelinks,bookmarks=false]{hyperref}
\newcommand{\ie}{i.e.\ }
\newcommand{\cf}{cf.\ }
\newcommand{\resp}{respectively\ }
\newcommand{\Subsection}{\subsection}
\providecommand{\nobreakdash}{\nobreak\hbox{-}}
\setlength{\emergencystretch}{2em}
\multlinegap0pt
\usepackage{amssymb}
\usepackage{enumerate}
\usepackage{tikz}
\newtheorem{theorem}{Theorem}
\title{SOME MONOCHROMATIC PATTERNS IN NATURAL NUMBERS}
\author{Arpita Ghosh, Surojit Ghosh}
\date{Source: arXiv:2411.14066v3 (CC BY 4.0)}
\begin{document}
\maketitle

\noindent\emph{Source and licence.} SOME MONOCHROMATIC PATTERNS IN NATURAL NUMBERS. Version arXiv:2411.14066v3 (CC BY 4.0), distributed by arXiv under the Creative Commons Attribution 4.0 International licence, http://creativecommons.org/licenses/by/4.0/, as recorded in the arXiv licence metadata for this exact version and shown on the abstract page https://arxiv.org/abs/2411.14066v3. Full licence notice: \texttt{licenses/arxiv-2411.14066-CC-BY-4.0.txt}.\par\smallskip

\noindent\textit{Editorial scope.} Counted unit: the article's theorem theorem (source0.tex lines 567--573). The article's complete original proof of it is delivered with it (source0.tex lines 575--600, one \texttt{proof} environment of 26 lines). No mathematics is written, repaired or re-derived: the statement and the proof are the article's own lines, in the article's own notation, and no retained line is substituted or re-flowed.  No further source-local context is needed: every cross-reference of every retained line resolves inside the two delivered spans.  Nothing was omitted from any retained span, so the package carries no omission record.  No retained line contains a citation, so the wrapper carries no bibliography and no bibliography entry.  No retained line contains a figure environment, an image-inclusion command or drawing code, so no image is delivered and the package declares no asset dependency.  Editorial formatting only, in addition: an AMS \texttt{amsart} preamble that restates the article's own declarations and macros used by the retained lines, namely the environments \texttt{theorem} on separate counters, together with the standard packages those lines need.  The retained environments print in this document as Theorem 1 in order of appearance, whereas the article prints the same items with the numbering of its own full document.  The complete native source of the article as distributed in the arXiv e-print for this exact version is bundled unchanged as \texttt{sources/168870/source0.tex}.
\begin{theorem}
Let \( d, k, \ell \in \mathbb{N} \) and \( \{F_1, \dots, F_\ell\} \subset \mathcal{P}_f(\mathbb{N}) \) with \( F_i = \{a_{i1}, \dots, a_{id}\} \) for all \( i \in \{1, \dots, \ell\} \). Then for any \( k \)-coloring of \( \mathbb{N} \), there exist \( b, c \in \mathbb{N} \) such that the set 
\[
\left\{\left|\Sigma_{<s_b s_{a_{i1}}^{c} \cdots s_{a_{id}}^{c^d}}\right| : i = 1, \dots, \ell \right\}
\]
is monochromatic, where \( \Sigma \) is the set of sums of two squares.
\end{theorem}

\begin{proof}
Let \( q, k, d \in \mathbb{N} \) be as in the statement of the polynomial Hales-Jewett theorem. By the theorem, there exists a natural number \( N = N(q, k, d) \). Suppose \( \chi : (\mathbb{N}, \ast_f) \to [k] \) is a \( k \)-coloring of \( \mathbb{N} \), and let \( m : X(q, N, d) \to (\mathbb{N}, \ast_f) \) be the canonical map defined by
\[
m((\mathbf{b}_1, \dots, \mathbf{b}_d)) = \overset{d}{\underset{j=1}{\ast_f}} \left( \underset{(i_1, \dots, i_j) \in [N]^j}{\ast_f} \mathbf{b}_j((i_1, \dots, i_j)) \right).
\]
The composite \( \chi \circ m \) is a \( k \)-coloring of \( X(q, N, d) \). By the polynomial Hales-Jewett theorem, there exist \( a \in X(q, N, d) \) and \( \gamma \subseteq [N] \) such that the set
\[
\{ a \oplus x_1 \gamma \oplus \cdots \oplus x_d \gamma^d : (x_1, \dots, x_d) \in [q]^d \}
\]
is monochromatic. Therefore, the image
\[
m\left(\{ a \oplus x_1 \gamma \oplus \cdots \oplus x_d \gamma^d : (x_1, \dots, x_d) \in [q]^d \}\right)
\]
is monochromatic for the coloring \( \chi \) of \( \mathbb{N} \). Note that
\[
m(a \oplus x_1 \gamma \oplus \cdots \oplus x_d \gamma^d) = \overset{d}{\underset{j=1}{\ast_f}} \left( \underset{(i_1, \dots, i_j) \in \gamma^j}{\ast_f} \mathbf{b}_j((i_1, \dots, i_j)) \right) \ast_f \left( \overset{d}{\underset{j=1}{\ast_f}} \left( \underset{(i_1, \dots, i_j) \in [N]^j \setminus \gamma^j}{\ast_f} \mathbf{b}_j((i_1, \dots, i_j)) \right) \right).
\]
This can be rewritten as
\[
x_1^{(c)} \ast_f \cdots \ast_f x_d^{(c^d)} \ast_f b = |\Sigma_{<  s_b s_{x_1}^{c} \cdots s_{x_d}^{c^d}}|,
\]
where \( c = |\gamma| \) and
\[
b = \overset{d}{\underset{j=1}{\ast_f}} \left( \underset{(i_1, \dots, i_j) \in [N]^j \setminus \gamma^j}{\ast_f} \mathbf{a}_j((i_1, \dots, i_j)) \right).
\]
\end{proof}

\end{document}
